\appendix
\chapter{Metodo dei minimi quadrati per il calcolo stabile della derivata}
\label{chap:derivata_regressione_lineare}
Di seguito viene illustrato il procedimento matematico effettuato per ottenere la formula utilizzata nel controllore PID per
il calcolo del contributo derivativo attraverso una regressione lineare sui dati contenuti nel buffer.

Il calcolo della derivata di una grandezza fisica misurata da un sensore risulta particolarmente delicato in quanto è fortemente
soggetto al rumore delle misurazioni. Il semplice calcolo del rapporto incrementale tra due campioni consecutivi può portare
a rapide variazioni della derivata che potrebbero indurre problemi di stabilità del controllore. Per mitigare le fluttuazioni
istantanee si è scelto di utilizzare il metodo dei minimi quadrati ed effettuare una regressione lineare su un set più ampio
di dati in modo da compensare le fluttuazioni istantanee e ottenere una stima più stabile della derivata.

\subsection*{Notazioni e convenzioni utilizzate}
Nei calcoli successivi si adottano le seguenti notazioni:
\begin{itemize}
	\item \(y_i\) : misurazione del sensore all'istante di tempo \(x_i\)
	\item \(a_0 + a_1 \, x_i\) : retta di regressione lineare, dove \(a_0\) è intercetta e \(a_1\) è il coefficiente angolare
	della retta di regressione lineare, che rappresenta la derivata da calcolare
\end{itemize}

Per calcolare la derivata ad un certo istante di tempo \(x_i\), si considerano solo gli ultimi \(N\) campioni misurati dal
sensore. Scelto \(x_k\) l'istante di tempo associato all'ultima misurazione \(y_j\), l'indice \(i\) dei campioni considerati
varia in \([j \!-\! (N \!-\! 1), \, j]\) con \(j \in \mathbb{N}\). Siccome gli indici delle sommatorie sono variabili mute,
nei successivi passaggi si considera senza perdita di generalità \(j = N \!-\! 1\) in modo che \(i\) varia da \(0\) a \(N-1\).

\subsection*{Procedimento matematico}
L'errore \(e_i\) tra la misurazione \(y_i\) e la retta di regressione lineare \(a_0 + a_1 \, x_i\) viene definito come:
\begin{equation}
	e_i = y_i - (a_0 + a_1 \, x_i)
\end{equation}

Si definisce quindi l'errore quadratico totale \(S(a_0, a_1)\) come segue:
\begin{equation}
	S(a_0, a_1) \; = \; \sum_{i = 0}^{N-1} e_i^2 \; = \; \sum_{i = 0}^{N-1} (y_i - (a_0 + a_1 \, x_i))^2
\end{equation}

I parametri \(a_0\) e \(a_1\) della retta di regressione lineare che minimizzano l'errore quadratico totale \(S(a_0, a_1)\)
si ottengono ponendo a zero le derivate parziali di \(S(a_0, a_1)\) rispetto a \(a_0\) e \(a_1\) e risolvendo il sistema
lineare di due equazioni a due incognite \(a_0\) e \(a_1\) così ottenuto.
\begin{align}
	\frac{\partial S(a_0, a_1)}{\partial a_0} &= -2 \sum_{i = 0}^{N-1} (y_i - (a_0 + a_1 \, x_i)) = \\
	&= -2 \left( \sum_{i = 0}^{N-1} y_i - \sum_{i=0}^{N-1} a_0 - \sum_{i=0}^{N-1} a_1 \, x_i \right) = \\
	&= -2 \left( \sum_{i = 0}^{N-1} y_i - a_0 \sum_{i=0}^{N-1} 1 - a_1 \sum_{i=0}^{N-1} x_i \right) = \\
	&= -2 \left( N \bar{y} - a_0 N - a_1 N \bar{x} \right) = \\
	&= -2 N \left( \bar{y} - a_0 - a_1 \bar{x} \right)
\end{align}
\begin{align}
	\frac{\partial S(a_0, a_1)}{\partial a_1} &= -2 \sum_{i = 0}^{N-1} x_i (y_i - (a_0 + a_1 \, x_i)) = \\
	&= -2 \left( \sum_{i = 0}^{N-1} x_i y_i - \sum_{i=0}^{N-1} a_0 x_i - \sum_{i=0}^{N-1} a_1 {x_i}^2 \right) = \\
	&= -2 \left( \sum_{i = 0}^{N-1} x_i y_i - a_0 \sum_{i=0}^{N-1} x_i - a_1 \sum_{i=0}^{N-1} {x_i}^2 \right) = \\
	&= -2 \left( \sum_{i = 0}^{N-1} x_i y_i - a_0 N \bar{x} - a_1 \sum_{i=0}^{N-1} {x_i}^2 \right)
\end{align}

\newpage

Si procede quindi a risolvere il sistema di equazioni lineari così ottenuto per sostituzione:
\begin{equation}
	\begin{cases}
		\frac{\partial S(a_0, a_1)}{\partial a_0} = 0 \\
		\frac{\partial S(a_0, a_1)}{\partial a_1} = 0
	\end{cases}
	\quad \Rightarrow \quad
	\begin{cases}
		2 N \left( \bar{y} - a_0 - a_1 \bar{x} \right) = 0 \\
		-2 \left( \sum_{i = 0}^{N-1} x_i y_i - a_0 N \bar{x} - a_1 \sum_{i=0}^{N-1} {x_i}^2 \right) = 0
	\end{cases}
\end{equation}

\begin{equation}
	-2 N \left( \bar{y} - a_0 - a_1 \bar{x} \right) = 0 \quad \Rightarrow \quad a_0 = \bar{y} - a_1 \bar{x}
\end{equation}
\begin{align}
	-2 \left( \sum_{i = 0}^{N-1} x_i y_i - a_0 N \bar{x} - a_1 \sum_{i=0}^{N-1} {x_i}^2 \right) = 0 \\
	\sum_{i = 0}^{N-1} x_i y_i - (\bar{y} - a_1 \bar{x}) N \bar{x} - a_1 \sum_{i=0}^{N-1} {x_i}^2 = 0 \\
	\sum_{i = 0}^{N-1} x_i y_i - N \bar{y} \bar{x} + a_1 N {\bar{x}}^2 - a_1 \sum_{i=0}^{N-1} {x_i}^2 = 0 \\
	\sum_{i = 0}^{N-1} x_i y_i - N \bar{y} \bar{x} - a_1 \left( \sum_{i=0}^{N-1} {x_i}^2 - N {\bar{x}}^2 \right) = 0 \\
\end{align}
\begin{equation}
	a_1 = \frac{\displaystyle \;\; \sum_{i=0}^{N-1} x_i y_i - N \bar{y} \bar{x} \;\;}{\displaystyle \;\; \sum_{i=0}^{N-1} {x_i}^2 - N {\bar{x}}^2 \;\;} \label{eq:7_a1_base}
\end{equation}

Si assume che gli istanti di tempo \(x_i\) sono tutti equidistanti, ovvero che formano una sequenza continua dei multipli del
tempo di campionamento \(T_s\) a partire dall'istante iniziale \(k \, T_s\) dove \(k = j - (N - 1)\). I campioni \(x_i\) possono
essere quindi espressi in funzione dei due indici \(i\) e \(k\):
\begin{equation}
	x_i = (j - (N-1)) \, T_s + i \, T_s = (k + i) \, T_s \qquad \text{con} \; \bar{x} = \frac{1}{N} \sum_{i=0}^{N-1} x_i = \frac{T_s}{N} \sum_{i=0}^{N-1} (k + i)\label{eq:7_sub}
\end{equation}

È possibile quindi riscrivere la formula \eqref{eq:7_a1_base} con la sostituzione \eqref{eq:7_sub} individuata sopra:
\begin{equation}
	a_1 = \frac{\displaystyle \;\; T_s \sum_{i=0}^{N-1} (k + i) \, y_i - N \bar{y} \left(\frac{T_s}{N} \, \sum_{i=0}^{N-1} (k + i)\right) \;\;}{\displaystyle \;\; {T_s}^2 \sum_{i=0}^{N-1} (k + i)^2 - N \left(\frac{T_s}{N} \,\sum_{i=0}^{N-1} (k + i)\right)^2 \;\;}
\end{equation}

\newpage

Analizzando separatamente il numeratore e il denominatore della formula \eqref{eq:7_a1_base} si ottengono le seguenti
semplificazioni fino ad ottenere la formula finale utilizzata nell'implementazione del controllore PID:
\begin{align}
	\textit{Num.} &= T_s \sum_{i=0}^{N-1} (k + i) \, y_i - N \bar{y} \left(\frac{T_s}{N} \, \sum_{i=0}^{N-1} (k + i)\right) = \\
	&= T_s \left(k \sum_{i=0}^{N-1} y_i + \sum_{i=0}^{N-1} i \cdot y_i\right) - N \bar{y} \cdot \frac{T_s}{N} \cdot \left( \sum_{i=0}^{N-1} i + \sum_{i=0}^{N-1} k \right) = \\
	&= T_s \left( k N \bar{y} + \sum_{i=0}^{N-1} i \cdot y_i - \bar{y} \left(\frac{N(N-1)}{2} + kN\right) \right) = \\
	&= T_s \left( k N \bar{y} + \sum_{i=0}^{N-1} i \cdot y_i - \frac{N(N-1)}{2} \cdot \bar{y} - k N \bar{y} \right) = \\
	&= T_s \left( \sum_{i=0}^{N-1} i \cdot y_i - \frac{N(N-1)}{2} \cdot \bar{y} \right) = \\
	&= T_s \left( \sum_{i=0}^{N-1} i \cdot y_i - \frac{N-1}{2} \cdot \sum_{i=0}^{N-1} y_i \right) \label{eq:7_num}
\end{align}
\begin{align}
	\textit{Den.} &= {T_s}^2 \sum_{i=0}^{N-1} (k + i)^2 - N \left(\frac{T_s}{N} \,\sum_{i=0}^{N-1} (k + i)\right)^2 = \\
	&= {T_s}^2 \left( \sum_{i=0}^{N-1} (k^2 + 2ik + i^2) - \frac{1}{N} \left( \sum_{i=0}^{N-1} (k + i) \right)^2 \right) = \\
	&= {T_s}^2 \left( k^2 N + 2k \sum_{i=0}^{N-1} i + \sum_{i=0}^{N-1} i^2 - \frac{1}{N} \left( k N + \sum_{i=0}^{N-1} i \right)^2 \right) = \\
	&= {T_s}^2 \left( k^2 N \!+\! 2k \frac{N(N\!\!-\!\!1)}{2} \!+\! \frac{N(N\!\!-\!\!1)(2N\!\!-\!\!1)}{6} - \frac{1}{N} \left( k N \!+\! \frac{N(N\!\!-\!\!1)}{2} \right)^2 \right) = \\
	\begin{split}
		&= {T_s}^2 \bigg( k^2 N + k N^2 - kN + \frac{2N^3 - 3 N^2 + N}{6} + \\[-5pt]
		&\qquad\qquad\qquad\qquad -\frac{1}{N} \cdot \bigg( k^2 N^2 + kN^3 - kN^2 + \frac{N^4 - 2N^3 + N^2}{4} \bigg) \bigg) = \\
	\end{split} \\
	\begin{split}
		&= \frac{{T_s}^2}{12} \cdot \big( 12k^2 N + 12k N^2 - 12kN + 4N^3 - 6N^2 + 2N + \\
		&\qquad\qquad\qquad\qquad + 12k^2 N - 12k N^2 + 12kN - 3N^3 + 6N^2 - 3N \big) =
	\end{split} \\
	&= \frac{{T_s}^2}{12} \left( N^3 - N \right) = \frac{{T_s}^2}{12} \cdot N (N-1)(N+1) \label{eq:7_den}
\end{align}


Unendo i due termini \eqref{eq:7_num} e \eqref{eq:7_den} così ottenuti, si arriva alla formula finale utilizzata
nell'implementazione del controllore PID.
\begin{align}
	a_1 &= \frac{\displaystyle \;\; T_s \sum_{i=0}^{N-1} (k + i) \, y_i - N \bar{y} \left(\frac{T_s}{N} \, \sum_{i=0}^{N-1} (k + i)\right) \;\;}{\displaystyle \;\; {T_s}^2 \sum_{i=0}^{N-1} (k + i)^2 - N \left(\frac{T_s}{N} \,\sum_{i=0}^{N-1} (k + i)\right)^2 \;\;} = \\[5pt]
	&= \frac{\displaystyle \;\; T_s \left( \sum_{i=0}^{N-1} i \cdot y_i - \frac{N-1}{2} \cdot \sum_{i=0}^{N-1} y_i \right) \;\;}{\displaystyle \;\; \frac{{T_s}^2}{12} \cdot N (N-1)(N+1) \;\;} = \\[5pt]
	&= \frac{12}{T_s \cdot N(N-1)(N+1)} \cdot \left( \sum_{i=0}^{N-1} i \cdot y_i - \frac{N-1}{2} \cdot \sum_{i=0}^{N-1} y_i \right) = \\
	&= \frac{12}{T_s \cdot N(N^2-1)} \cdot \sum_{i=0}^{N-1} i \cdot y_i - \frac{6}{T_s \cdot N(N+1)} \cdot \sum_{i=0}^{N-1} y_i
\end{align}


% -------------------------------------------------------------------------------------------------
% ------------- codice con verbatim se è richiesto il formato PDF/A-1a (accessibile) --------------
% -------------------------------------------------------------------------------------------------
%
% Per utilizzare correttamente il sistema di tagging dei listing richiesto dallo standard
% PDF/A-1a (accessibile) non è possibile utilizzare il comando \captionof{}{} al di fuori
% di ambienti flottanti. Per questo motivo, si riporta di seguito il procedimento da seguire
% in 5 fasi principali:
%
%  1. incremento del contatore dei listing
%     \noindent\refstepcounter{lstlisting}
%
%  2. creazione della reference per il listing
%	  \label{lst:label_del_listing}
%
%  3. stampa della didascalia del listing
%     \noindent\text{Codice \thelstlisting: Didascalia del listing}
%
%  4. aggiunta del listing all'indice dei listing
%     \addcontentsline{lol}{lstlisting}{\protect\numberline{\thelstlisting}Didascalia del listing}
%
%  5. importazione del file di codice con verbatim aggiungendo le linee orizzontali sopra e sotto il codice
%     per una migliore leggibilità e con dimensione del font ridotta
%     \noindent\rule{\textwidth}{0.4pt}
%     {\scriptsize \verbatiminput{percorso_del_file}}
%     \noindent\rule{\textwidth}{0.4pt}
%

%\chapter{Codice del firmware}
%\section{File di configurazione per PlatformIO in VSCode}
%
%% etichetta file di configurazione dell'ambiente PlatformIO in VSCode
%\noindent\refstepcounter{lstlisting}
%\label{lst:platformio_ini}
%\noindent\text{Codice \thelstlisting: File di configurazione dell'ambiente PlatformIO in VSCode}
%\addcontentsline{lol}{lstlisting}{\protect\numberline{\thelstlisting}File di configurazione dell'ambiente PlatformIO in VSCode}
%
%% codice del file di configurazione dell'ambiente PlatformIO in VSCode
%\noindent\rule{\textwidth}{0.4pt}
%{\tiny \verbatiminput{../programmi/firmware/platformio.ini}}
%\noindent\rule{\textwidth}{0.4pt}
%
%\section{File di configurazione del firmware}
%
%% etichetta file header con i parametri di configurazione
%\noindent\refstepcounter{lstlisting}
%\label{lst:config_h}
%\noindent\text{Codice \thelstlisting: File header con i parametri di configurazione}
%\addcontentsline{lol}{lstlisting}{\protect\numberline{\thelstlisting}File header con i parametri di configurazione}
%
%% codice del file header con i parametri di configurazione
%\noindent\rule{\textwidth}{0.4pt}
%{\tiny \verbatiminput{../programmi/firmware/lib/config/config.h}}
%\noindent\rule{\textwidth}{0.4pt}
%
%% etichetta file sorgente con i parametri di configurazione
%\noindent\refstepcounter{lstlisting}
%\label{lst:config_cpp}
%\noindent\text{Codice \thelstlisting: File sorgente con i parametri di configurazione}
%\addcontentsline{lol}{lstlisting}{\protect\numberline{\thelstlisting}File sorgente con i parametri di configurazione}
%
%% codice del file sorgente con i parametri di configurazione
%\noindent\rule{\textwidth}{0.4pt}
%{\tiny \verbatiminput{../programmi/firmware/lib/config/config.cpp}}
%\noindent\rule{\textwidth}{0.4pt}
%
%\newpage
%
%\section{File del modulo con le funzioni di utilità}
%
%% etichetta file header con le funzioni di utilità
%\noindent\refstepcounter{lstlisting}
%\label{lst:utils_h}
%\noindent\text{Codice \thelstlisting: File header con le funzioni di utilità}
%\addcontentsline{lol}{lstlisting}{\protect\numberline{\thelstlisting}File header con le funzioni di utilità}
%
%% codice del file header con le funzioni di utilità
%\noindent\rule{\textwidth}{0.4pt}
%{\tiny \verbatiminput{../programmi/firmware/lib/utils/utils.h}}
%\noindent\rule{\textwidth}{0.4pt}
%
%% etichetta file sorgente con le funzioni di utilità
%\noindent\refstepcounter{lstlisting}
%\label{lst:utils_cpp}
%\noindent\text{Codice \thelstlisting: File sorgente con le funzioni di utilità}
%\addcontentsline{lol}{lstlisting}{\protect\numberline{\thelstlisting}File sorgente con le funzioni di utilità}
%
%% codice del file sorgente con le funzioni di utilità
%\noindent\rule{\textwidth}{0.4pt}
%{\tiny \verbatiminput{../programmi/firmware/lib/utils/utils.cpp}}
%\noindent\rule{\textwidth}{0.4pt}
%
%\newpage
%
%\section{File del modulo per la gestione degli errori}
%
%% etichetta file header con le funzioni per la gestione degli errori
%\noindent\refstepcounter{lstlisting}
%\label{lst:error_h}
%\noindent\text{Codice \thelstlisting: File header con le funzioni per la gestione degli errori}
%\addcontentsline{lol}{lstlisting}{\protect\numberline{\thelstlisting}File header con le funzioni per la gestione degli errori}
%
%% codice del file header con le funzioni per la gestione degli errori
%\noindent\rule{\textwidth}{0.4pt}
%{\tiny \verbatiminput{../programmi/firmware/lib/error/error.h}}
%\noindent\rule{\textwidth}{0.4pt}
%
%% etichetta file sorgente con le funzioni per la gestione degli errori
%\noindent\refstepcounter{lstlisting}
%\label{lst:error_cpp}
%\noindent\text{Codice \thelstlisting: File sorgente con le funzioni per la gestione degli errori}
%\addcontentsline{lol}{lstlisting}{\protect\numberline{\thelstlisting}File sorgente con le funzioni per la gestione degli errori}
%
%% codice del file sorgente con le funzioni per la gestione degli errori
%\noindent\rule{\textwidth}{0.4pt}
%{\tiny \verbatiminput{../programmi/firmware/lib/error/error.cpp}}
%\noindent\rule{\textwidth}{0.4pt}
%
%\newpage
%
%\section{File del modulo per il controllo attuatori}
%
%% etichetta file header con le funzioni per il controllo degli attuatori
%\noindent\refstepcounter{lstlisting}
%\label{lst:actuators_h}
%\noindent\text{Codice \thelstlisting: File header con le funzioni per il controllo degli attuatori}
%\addcontentsline{lol}{lstlisting}{\protect\numberline{\thelstlisting}File header con le funzioni per il controllo degli attuatori}
%
%% codice del file header con le funzioni per il controllo degli attuatori
%\noindent\rule{\textwidth}{0.4pt}
%{\tiny \verbatiminput{../programmi/firmware/lib/actuators/actuators.h}}
%\noindent\rule{\textwidth}{0.4pt}
%
%% etichetta file sorgente con le funzioni per il controllo degli attuatori
%\noindent\refstepcounter{lstlisting}
%\label{lst:actuators_cpp}
%\noindent\text{Codice \thelstlisting: File sorgente con le funzioni per il controllo degli attuatori}
%\addcontentsline{lol}{lstlisting}{\protect\numberline{\thelstlisting}File sorgente con le funzioni per il controllo degli attuatori}
%
%% codice del file sorgente con le funzioni per il controllo degli attuatori
%\noindent\rule{\textwidth}{0.4pt}
%{\tiny \verbatiminput{../programmi/firmware/lib/actuators/actuators.cpp}}
%\noindent\rule{\textwidth}{0.4pt}
%
%\newpage
%
%\section{File del modulo per la gestione del display lcd}
%
%% etichetta file header con le funzioni per la gestione del display lcd
%\noindent\refstepcounter{lstlisting}
%\label{lst:lcd_h}
%\noindent\text{Codice \thelstlisting: File header con le funzioni per la gestione del display lcd}
%\addcontentsline{lol}{lstlisting}{\protect\numberline{\thelstlisting}File header con le funzioni per la gestione del display lcd}
%
%% codice del file header con le funzioni per la gestione del display lcd
%\noindent\rule{\textwidth}{0.4pt}
%{\tiny \verbatiminput{../programmi/firmware/lib/lcd/lcd.h}}
%\noindent\rule{\textwidth}{0.4pt}
%
%% etichetta file sorgente con le funzioni per la gestione del display lcd
%\noindent\refstepcounter{lstlisting}
%\label{lst:lcd_cpp}
%\noindent\text{Codice \thelstlisting: File sorgente con le funzioni per la gestione del display lcd}
%\addcontentsline{lol}{lstlisting}{\protect\numberline{\thelstlisting}File sorgente con le funzioni per la gestione del display lcd}
%
%% codice del file sorgente con le funzioni per la gestione del display lcd
%\noindent\rule{\textwidth}{0.4pt}
%{\tiny \verbatiminput{../programmi/firmware/lib/lcd/lcd.cpp}}
%\noindent\rule{\textwidth}{0.4pt}
%
%\newpage
%
%\section{File del modulo per la gestione del sensore di temperatura e umidità SHT20}
%
%% etichetta file header con le funzioni per la gestione del sensore SHT20
%\noindent\refstepcounter{lstlisting}
%\label{lst:sht20_h}
%\noindent\text{Codice \thelstlisting: File header con le funzioni per la gestione del sensore SHT20}
%\addcontentsline{lol}{lstlisting}{\protect\numberline{\thelstlisting}File header con le funzioni per la gestione del sensore SHT20}
%
%% codice del file header con le funzioni per la gestione del sensore SHT20
%\noindent\rule{\textwidth}{0.4pt}
%{\tiny \verbatiminput{../programmi/firmware/lib/sht20/sht20.h}}
%\noindent\rule{\textwidth}{0.4pt}
%
%% etichetta file sorgente con le funzioni per la gestione del sensore SHT20
%\noindent\refstepcounter{lstlisting}
%\label{lst:sht20_cpp}
%\noindent\text{Codice \thelstlisting: File sorgente con le funzioni per la gestione del sensore SHT20}
%\addcontentsline{lol}{lstlisting}{\protect\numberline{\thelstlisting}File sorgente con le funzioni per la gestione del sensore SHT20}
%
%% codice del file sorgente con le funzioni per la gestione del sensore SHT20
%\noindent\rule{\textwidth}{0.4pt}
%{\tiny \verbatiminput{../programmi/firmware/lib/sht20/sht20.cpp}}
%\noindent\rule{\textwidth}{0.4pt}
%
%\newpage
%
%\section{File del modulo per la gestione delle comunicazioni seriali}
%
%% etichetta file header con le funzioni per la gestione delle comunicazioni seriali
%\noindent\refstepcounter{lstlisting}
%\label{lst:serial_h}
%\noindent\text{Codice \thelstlisting: File header con le funzioni per la gestione delle comunicazioni seriali}
%\addcontentsline{lol}{lstlisting}{\protect\numberline{\thelstlisting}File header con le funzioni per la gestione delle comunicazioni seriali}
%
%% codice del file header con le funzioni per la gestione delle comunicazioni seriali
%\noindent\rule{\textwidth}{0.4pt}
%{\tiny \verbatiminput{../programmi/firmware/lib/serial/serial.h}}
%\noindent\rule{\textwidth}{0.4pt}
%
%% etichetta file sorgente con le funzioni per la gestione delle comunicazioni seriali
%\noindent\refstepcounter{lstlisting}
%\label{lst:serial_cpp}
%\noindent\text{Codice \thelstlisting: File sorgente con le funzioni per la gestione delle comunicazioni seriali}
%\addcontentsline{lol}{lstlisting}{\protect\numberline{\thelstlisting}File sorgente con le funzioni per la gestione delle comunicazioni seriali}
%
%% codice del file sorgente con le funzioni per la gestione delle comunicazioni seriali
%\noindent\rule{\textwidth}{0.4pt}
%{\tiny \verbatiminput{../programmi/firmware/lib/serial/serial.cpp}}
%\noindent\rule{\textwidth}{0.4pt}
%
%\newpage
%
%\section{File del modulo per la gestione del controllore PID}
%
%% etichetta file header con le funzioni per la gestione del controllore PID
%\noindent\refstepcounter{lstlisting}
%\label{lst:pid_h}
%\noindent\text{Codice \thelstlisting: File header con le funzioni per la gestione del controllore PID}
%\addcontentsline{lol}{lstlisting}{\protect\numberline{\thelstlisting}File header con le funzioni per la gestione del controllore PID}
%
%% codice del file header con le funzioni per la gestione del controllore PID
%\noindent\rule{\textwidth}{0.4pt}
%{\tiny \verbatiminput{../programmi/firmware/lib/pid/pid.h}}
%\noindent\rule{\textwidth}{0.4pt}
%
%\newpage
%
%\section{File principale del firmware}
%
%% etichetta file principale del firmware
%\noindent\refstepcounter{lstlisting}
%\label{lst:main_cpp}
%\noindent\text{Codice \thelstlisting: File principale del firmware}
%\addcontentsline{lol}{lstlisting}{\protect\numberline{\thelstlisting}File principale del firmware}
%
%% codice del file principale del firmware
%\noindent\rule{\textwidth}{0.4pt}
%{\tiny \verbatiminput{../programmi/firmware/src/main.cpp}}
%\noindent\rule{\textwidth}{0.4pt}
%
%\chapter{Programmi in MATLAB}
%\section{Gestione della comunicazione seriale}
%
%% etichetta funzione MATLAB per la gestione della comunicazione seriale
%\noindent\refstepcounter{lstlisting}
%\label{lst:serial_read}
%\noindent\text{Codice \thelstlisting: Funzione MATLAB che si occupa della lettura dei messaggi dalla porta seriale, del
%relativo parsing e salvataggio in un file csv e aggiornare i grafici di andamento delle variabili in tempo reale.}
%\addcontentsline{lol}{lstlisting}{\protect\numberline{\thelstlisting} Funzione MATLAB per la gestione della comunicazione seriale}
%
%% codice della funzione MATLAB per la gestione della comunicazione seriale
%\noindent\rule{\textwidth}{0.4pt}
%{\tiny \verbatiminput{../matlab/serial_read.m}}
%\noindent\rule{\textwidth}{0.4pt}
%
%\newpage
%
%\section{Visualizzazione dei grafici dei dati csv}
%
%% etichetta funzione MATLAB per la visualizzazione dei grafici dei dati csv
%\noindent\refstepcounter{lstlisting}
%\label{lst:plot_csv}
%\noindent\text{Codice \thelstlisting: Funzione MATLAB per la stampa dei grafici dell'andamento delle grandezze fisiche in
%funzione del tempo, utilizzando i dati raccolti nei file csv.}
%\addcontentsline{lol}{lstlisting}{\protect\numberline{\thelstlisting}Funzione MATLAB per visualizzare i grafici dei dati csv}
%
%% codice della funzione MATLAB per la visualizzazione dei grafici dei dati csv
%\noindent\rule{\textwidth}{0.4pt}
%{\tiny \verbatiminput{../matlab/plot_csv.m}}
%\noindent\rule{\textwidth}{0.4pt}
%
%\newpage
%
%\section{Analisi dei dati acquisiti durante i test}
%
%% etichetta funzione MATLAB per l'analisi dati del primo test
%\noindent\refstepcounter{lstlisting}
%\label{lst:test1_analysis}
%\noindent\text{Codice \thelstlisting: Funzione MATLAB per il calcolo del periodo di warm-up a partire dai dati raccolti nei
%file csv durante il primo test previsto dalla normativa.}
%\addcontentsline{lol}{lstlisting}{\protect\numberline{\thelstlisting}Funzione MATLAB per l'analisi dati del primo test}
%
%% codice della funzione MATLAB per l'analisi dati del primo test
%\noindent\rule{\textwidth}{0.4pt}
%{\tiny \verbatiminput{../matlab/test1_analysis.m}}
%\noindent\rule{\textwidth}{0.4pt}
%
%% etichetta funzione MATLAB per l'analisi dati del secondo test
%\noindent\refstepcounter{lstlisting}
%\label{lst:test2_analysis}
%\noindent\text{Codice \thelstlisting: Funzione MATLAB per il calcolo delle metriche di prestazione (overshoot, settling time
%ed errori a regime) con i dati raccolti nei file csv durante il secondo test previsto dalla normativa.}
%\addcontentsline{lol}{lstlisting}{\protect\numberline{\thelstlisting}Funzione MATLAB per l'analisi dati del secondo test}
%
%% codice della funzione MATLAB per l'analisi dati del secondo test
%\noindent\rule{\textwidth}{0.4pt}
%{\tiny \verbatiminput{../matlab/test2_analysis.m}}
%\noindent\rule{\textwidth}{0.4pt}


% -------------------------------------------------------------------------------------------------
% ---------------- codice con minted se è sufficiente il formato PDF/A-2b (basic) -----------------
% -------------------------------------------------------------------------------------------------

\chapter{Codice del firmware}
\section{File di configurazione per PlatformIO in VSCode}

% etichetta file di configurazione dell'ambiente PlatformIO in VSCode
\noindent\refstepcounter{listing}
\label{lst:platformio_ini}
\noindent\text{Codice \thelisting: File di configurazione dell'ambiente PlatformIO in VSCode}
\addcontentsline{lol}{listing}{\protect\numberline{\thelisting}File di configurazione dell'ambiente PlatformIO in VSCode}
\inputminted{ini}{../programmi/firmware/platformio.ini}

\section{File di configurazione del firmware}

% etichetta file header con i parametri di configurazione
\noindent\refstepcounter{listing}
\label{lst:config_h}
\noindent\text{Codice \thelisting: File header con i parametri di configurazione}
\addcontentsline{lol}{listing}{\protect\numberline{\thelisting}File header con i parametri di configurazione}
\inputminted{c++}{../programmi/firmware/lib/config/config.h}

% etichetta file sorgente con i parametri di configurazione
\noindent\refstepcounter{listing}
\label{lst:config_cpp}
\noindent\text{Codice \thelisting: File sorgente con i parametri di configurazione}
\addcontentsline{lol}{listing}{\protect\numberline{\thelisting}File sorgente con i parametri di configurazione}
\inputminted{c++}{../programmi/firmware/lib/config/config.cpp}

\newpage

\section{File del modulo con le funzioni di utilità}

% etichetta file header con le funzioni di utilità
\noindent\refstepcounter{listing}
\label{lst:utils_h}
\noindent\text{Codice \thelisting: File header con le funzioni di utilità}
\addcontentsline{lol}{listing}{\protect\numberline{\thelisting}File header con le funzioni di utilità}
\inputminted{c++}{../programmi/firmware/lib/utils/utils.h}

% etichetta file sorgente con le funzioni di utilità
\noindent\refstepcounter{listing}
\label{lst:utils_cpp}
\noindent\text{Codice \thelisting: File sorgente con le funzioni di utilità}
\addcontentsline{lol}{listing}{\protect\numberline{\thelisting}File sorgente con le funzioni di utilità}
\inputminted{c++}{../programmi/firmware/lib/utils/utils.cpp}

\newpage

\section{File del modulo per la gestione degli errori}

% etichetta file header con le funzioni per la gestione degli errori
\noindent\refstepcounter{listing}
\label{lst:error_h}
\noindent\text{Codice \thelisting: File header con le funzioni per la gestione degli errori}
\addcontentsline{lol}{listing}{\protect\numberline{\thelisting}File header con le funzioni per la gestione degli errori}
\inputminted{c++}{../programmi/firmware/lib/error/error.h}

% etichetta file sorgente con le funzioni per la gestione degli errori
\noindent\refstepcounter{listing}
\label{lst:error_cpp}
\noindent\text{Codice \thelisting: File sorgente con le funzioni per la gestione degli errori}
\addcontentsline{lol}{listing}{\protect\numberline{\thelisting}File sorgente con le funzioni per la gestione degli errori}
\inputminted{c++}{../programmi/firmware/lib/error/error.cpp}

\newpage

\section{File del modulo per il controllo attuatori}

% etichetta file header con le funzioni per il controllo degli attuatori
\noindent\refstepcounter{listing}
\label{lst:actuators_h}
\noindent\text{Codice \thelisting: File header con le funzioni per il controllo degli attuatori}
\addcontentsline{lol}{listing}{\protect\numberline{\thelisting}File header con le funzioni per il controllo degli attuatori}
\inputminted{c++}{../programmi/firmware/lib/actuators/actuators.h}

% etichetta file sorgente con le funzioni per il controllo degli attuatori
\noindent\refstepcounter{listing}
\label{lst:actuators_cpp}
\noindent\text{Codice \thelisting: File sorgente con le funzioni per il controllo degli attuatori}
\addcontentsline{lol}{listing}{\protect\numberline{\thelisting}File sorgente con le funzioni per il controllo degli attuatori}
\inputminted{c++}{../programmi/firmware/lib/actuators/actuators.cpp}

\newpage

\section{File del modulo per la gestione del display lcd}

% etichetta file header con le funzioni per la gestione del display lcd
\noindent\refstepcounter{listing}
\label{lst:lcd_h}
\noindent\text{Codice \thelisting: File header con le funzioni per la gestione del display lcd}
\addcontentsline{lol}{listing}{\protect\numberline{\thelisting}File header con le funzioni per la gestione del display lcd}
\inputminted{c++}{../programmi/firmware/lib/lcd/lcd.h}

% etichetta file sorgente con le funzioni per la gestione del display lcd
\noindent\refstepcounter{listing}
\label{lst:lcd_cpp}
\noindent\text{Codice \thelisting: File sorgente con le funzioni per la gestione del display lcd}
\addcontentsline{lol}{listing}{\protect\numberline{\thelisting}File sorgente con le funzioni per la gestione del display lcd}
\inputminted{c++}{../programmi/firmware/lib/lcd/lcd.cpp}

\newpage

\section{File del modulo per la gestione del sensore di temperatura e umidità SHT20}

% etichetta file header con le funzioni per la gestione del sensore SHT20
\noindent\refstepcounter{listing}
\label{lst:sht20_h}
\noindent\text{Codice \thelisting: File header con le funzioni per la gestione del sensore SHT20}
\addcontentsline{lol}{listing}{\protect\numberline{\thelisting}File header con le funzioni per la gestione del sensore SHT20}
\inputminted{c++}{../programmi/firmware/lib/sht20/sht20.h}

% etichetta file sorgente con le funzioni per la gestione del sensore SHT20
\noindent\refstepcounter{listing}
\label{lst:sht20_cpp}
\noindent\text{Codice \thelisting: File sorgente con le funzioni per la gestione del sensore SHT20}
\addcontentsline{lol}{listing}{\protect\numberline{\thelisting}File sorgente con le funzioni per la gestione del sensore SHT20}
\inputminted{c++}{../programmi/firmware/lib/sht20/sht20.cpp}

\newpage

\section{File del modulo per la gestione delle comunicazioni seriali}

% etichetta file header con le funzioni per la gestione delle comunicazioni seriali
\noindent\refstepcounter{listing}
\label{lst:serial_h}
\noindent\text{Codice \thelisting: File header con le funzioni per la gestione delle comunicazioni seriali}
\addcontentsline{lol}{listing}{\protect\numberline{\thelisting}File header con le funzioni per la gestione delle comunicazioni seriali}
\inputminted{c++}{../programmi/firmware/lib/serial/serial.h}

% etichetta file sorgente con le funzioni per la gestione delle comunicazioni seriali
\noindent\refstepcounter{listing}
\label{lst:serial_cpp}
\noindent\text{Codice \thelisting: File sorgente con le funzioni per la gestione delle comunicazioni seriali}
\addcontentsline{lol}{listing}{\protect\numberline{\thelisting}File sorgente con le funzioni per la gestione delle comunicazioni seriali}
\inputminted{c++}{../programmi/firmware/lib/serial/serial.cpp}

\newpage

\section{File del modulo per la gestione del controllore PID}

% etichetta file header con le funzioni per la gestione del controllore PID
\noindent\refstepcounter{listing}
\label{lst:pid_h}
\noindent\text{Codice \thelisting: File header con le funzioni per la gestione del controllore PID}
\addcontentsline{lol}{listing}{\protect\numberline{\thelisting}File header con le funzioni per la gestione del controllore PID}
\inputminted{c++}{../programmi/firmware/lib/pid/pid.h}

\newpage

\section{File principale del firmware}

% etichetta file principale del firmware
\noindent\refstepcounter{listing}
\label{lst:main_cpp}
\noindent\text{Codice \thelisting: File principale del firmware}
\addcontentsline{lol}{listing}{\protect\numberline{\thelisting}File principale del firmware}
\inputminted{c++}{../programmi/firmware/src/main.cpp}

\chapter{Programmi in MATLAB}
\section{Gestione della comunicazione seriale}

% etichetta funzione MATLAB per la gestione della comunicazione seriale
\noindent\refstepcounter{listing}
\label{lst:serial_read}
\noindent
Codice \thelisting: Funzione MATLAB che si occupa della lettura dei messaggi dalla porta seriale, del relativo parsing
e salvataggio in un file csv e aggiornare i grafici di andamento delle variabili in tempo reale.
\addcontentsline{lol}{listing}{\protect\numberline{\thelisting} Funzione MATLAB per la gestione della comunicazione seriale}
\inputminted[tabsize=4]{matlab}{../matlab/serial_read.m}

\newpage

\section{Visualizzazione dei grafici dei dati csv}

% etichetta funzione MATLAB per la visualizzazione dei grafici dei dati csv
\noindent\refstepcounter{listing}
\label{lst:plot_csv}
\noindent
Codice \thelisting: Funzione MATLAB per la stampa dei grafici dell'andamento delle grandezze fisiche in funzione del tempo,
utilizzando i dati raccolti nei file csv.
\addcontentsline{lol}{listing}{\protect\numberline{\thelisting}Funzione MATLAB la visualizzazione dei grafici dei dati csv}
\inputminted[tabsize=4]{matlab}{../matlab/plot_csv.m}

\newpage

\section{Analisi dei dati acquisiti durante i test}

% etichetta funzione MATLAB per l'analisi dati del primo test
\noindent\refstepcounter{listing}
\label{lst:test1_analysis}
\noindent
Codice \thelisting: Funzione MATLAB per il calcolo del periodo di warm-up a partire dai dati raccolti nei file csv durante
il primo test previsto dalla normativa.
\addcontentsline{lol}{listing}{\protect\numberline{\thelisting}Funzione MATLAB per l'analisi dati del primo test}
\inputminted[tabsize=4]{matlab}{../matlab/test1_analysis.m}

% etichetta funzione MATLAB per l'analisi dati del secondo test
\noindent\refstepcounter{listing}
\label{lst:test2_analysis}
\noindent
Codice \thelisting: Funzione MATLAB per il calcolo delle metriche di prestazione (overshoot, settling time ed errori a
regime) con i dati raccolti nei file csv durante il secondo test previsto dalla normativa.
\addcontentsline{lol}{listing}{\protect\numberline{\thelisting}Funzione MATLAB per l'analisi dati del secondo test}
\inputminted[tabsize=4]{matlab}{../matlab/test2_analysis.m}
